\subsection{Homological dimension of polynomial rings}

Lemma 2. --- Let \(\rho: \mathbf{A} \to \mathbf{A}'\) be a ring homomorphism, M an A-module, \(\mathbf{M}'\) a \(\mathbf{A}'\) -module. If the A-module \(\mathbf{A}_d'\) is flat, we have \(\mathrm{dp}_{\mathbf{A}'}(\mathbf{A}' \otimes_{\mathbf{A}} \mathbf{M}) \leqslant \mathrm{dp}_{\mathbf{A}}(\mathbf{M})\) . If the A-module \(\mathbf{A}_s'\) is projective, we have \(\mathrm{dp}_{\mathbf{A}}(\mathbf{M}') \leqslant \mathrm{dp}_{\mathbf{A}'}(\mathbf{M}')\) .

The first assertion is clear if \(\mathrm{dp}_{\mathbf{A}}(\mathbf{M}) = \pm \infty\) ; if \(\mathrm{dp}_{\mathbf{A}}(\mathbf{M}) = n \in \mathbb{N}\) , there exists an exact sequence of A-modules

\[
0 \rightarrow \mathrm{P} _ {n} \rightarrow \mathrm{P} _ {n - 1} \rightarrow \dots \rightarrow \mathrm{P} _ {0} \rightarrow \mathrm{M} \rightarrow 0
\]

where the \(P_{i}\) are projective; the sequence of A'-modules

\[
0 \rightarrow \mathrm{A} ^ {\prime} \otimes_ {\mathrm{A}} \mathrm{P} _ {n} \rightarrow \mathrm{A} ^ {\prime} \otimes_ {\mathrm{A}} \mathrm{P} _ {n - 1} \rightarrow \dots \rightarrow \mathrm{A} ^ {\prime} \otimes_ {\mathrm{A}} \mathrm{P} _ {0} \rightarrow \mathrm{A} ^ {\prime} \otimes_ {\mathrm{A}} \mathrm{M} \rightarrow 0
\]

is exact, since \(\mathbf{A}_d'\) is flat, and the \(\mathbf{A}'\) -modules \(\mathbf{A}' \otimes_{\mathbf{A}} \mathrm{P}_i\) are projective (II, p.~89, cor.); hence \(\mathrm{dp}_{\mathbf{A}'}(\mathbf{A}' \otimes_{\mathbf{A}} \mathbf{M}) \leqslant n = \mathrm{dp}_{\mathbf{A}}(\mathbf{M})\) . The second assertion is clear if \(\mathrm{dp}_{\mathbf{A}'}(\mathbf{M}') = \pm \infty\) ; if \(\mathrm{dp}_{\mathbf{A}'}(\mathbf{M}') = m \in \mathbf{N}\) , there exists an exact sequence of \(\mathbf{A}'\) -modules

\[
0 \to \mathrm{P} _ {m} ^ {\prime} \to \mathrm{P} _ {m - 1} ^ {\prime} \to \dots \to \mathrm{P} _ {0} ^ {\prime} \to \mathrm{M} ^ {\prime} \to 0,
\]

where the \(P_{i}^{\prime}\) are projective; the sequence of underlying A-modules is exact. On the other hand, each \(P_{i}^{\prime}\) is a direct summand of an A'-submodule of a module \(A_{s}^{\prime(I)}\) , hence is a projective A-module; we therefore have \(\mathrm{dp}_{\mathbf{A}}(\mathbf{M}^{\prime}) \leqslant m = \mathrm{dp}_{\mathbf{A}^{\prime}}(\mathbf{M}^{\prime})\) .

Lemma 3. --- Suppose A is commutative and let M be an A{[}X{]}-module.

\begin{enumerate}
\def\labelenumi{\alph{enumi})}
\item
  We have \(\mathrm{dp}_{\mathbf{A}}(\mathbf{M}) \leqslant \mathrm{dp}_{\mathbf{A}[\mathbf{X}]}(\mathbf{M}) \leqslant \mathrm{dp}_{\mathbf{A}}(\mathbf{M}) + 1\) .
\item
  If the homothety \(X_{M}\) is injective, we have dp \(_{A}\) (M/XM) ≤ dp \(_{A[X]}\) (M).
\item
  If \(\mathbf{XM} = 0\) , we have \(\mathrm{dp}_{\mathbf{A}}(\mathbf{M}) + 1 = \mathrm{dp}_{\mathbf{A}[\mathbf{X}]}(\mathbf{M}).\)
\item
  We have an exact sequence of A{[}X{]}-modules (III, p.~106 and VII, § 5, no. 1)
\end{enumerate}

\[
0 \rightarrow \mathrm{A} [ \mathrm{X} ] \otimes_ {\mathrm{A}} \mathrm{M} \rightarrow \mathrm{A} [ \mathrm{X} ] \otimes_ {\mathrm{A}} \mathrm{M} \rightarrow \mathrm{M} \rightarrow 0;
\]

Assertion (a) then follows from X, p.~135, cor. 2 and from Lemma 2.

\begin{enumerate}
\def\labelenumi{\alph{enumi})}
\setcounter{enumi}{1}
\tightlist
\item
  If \(\mathrm{dp}_{\mathbf{A}[\mathbf{X}]}(\mathbf{M}) = \pm \infty\) the assertion is trivial. If \(\mathbf{M}\) is projective and nonzero, then the A-module \(\mathbf{M} / \mathbf{XM}\) is identified with \(\mathbf{A} \otimes_{\mathbf{A}[\mathbf{X}]} \mathbf{M}\) , hence is projective, and we have \(\mathrm{dp}_{\mathbf{A}}(\mathbf{M} / \mathbf{XM}) \leqslant 0 = \mathrm{dp}_{\mathbf{A}[\mathbf{X}]}(\mathbf{M})\) . Let us reason by induction on \(\mathrm{dp}_{\mathbf{A}[\mathbf{X}]}(\mathbf{M}) = n\) , assumed \(>0\) . Consider an exact sequence of \(\mathbf{A}[\mathbf{X}]\) -modules \(0 \to \mathbf{N} \to \mathbf{L} \to \mathbf{M} \to 0\) , where \(\mathbf{L}\) is a free \(\mathbf{A}[\mathbf{X}]\) -module; applying to the diagram
\end{enumerate}

\[
\begin{array}{c} 0 \longrightarrow \mathrm{N} \longrightarrow \mathrm{L} \longrightarrow \mathrm{M} \longrightarrow 0 \\ \mathrm {X_ {N}} \Big \downarrow \quad \mathrm {X_ {L}} \Big \downarrow \quad \mathrm {X_ {M}} \Big \downarrow \\ 0 \longrightarrow \mathrm{N} \longrightarrow \mathrm{L} \longrightarrow \mathrm{M} \longrightarrow 0 \end{array}
\]

prop. 2 of X, p.~4, we see that \(X_{N}\) is injective and that we have the exact sequence

\[
0 \rightarrow \mathrm{N} / \mathrm{XN} \rightarrow \mathrm{L} / \mathrm{XL} \rightarrow \mathrm{M} / \mathrm{XM} \rightarrow 0.
\]

Since L is free over A{[}X{]}, L/XL is free over A and we have

\[
\mathrm{dp} _ {\mathbf {A} [ \mathbf {X} ]} (\mathbf {N}) = n - 1, \quad \mathrm{dp} _ {\mathbf {A}} (\mathbf {M} / \mathbf {X M}) \leqslant 1 + \mathrm{dp} _ {\mathbf {A}} (\mathbf {N} / \mathbf {X N});
\]

as dp \(_{A}\) (N/XN) ≤ n - 1 by the induction hypothesis, we deduce dp \(_{A}\) (M/XM) ≤ n, which is what had to be proved.

\begin{enumerate}
\def\labelenumi{\alph{enumi})}
\setcounter{enumi}{2}
\tightlist
\item
  The assertion is trivial if \(\mathrm{dp}_{\mathbf{A}[\mathbf{X}]}(\mathbf{M}) = \pm \infty\) , and also if \(\mathrm{dp}_{\mathbf{A}[\mathbf{X}]}(\mathbf{M}) = 0\) (which is impossible since \(\mathbf{XM} = 0\) ). We may therefore assume \(\mathrm{dp}_{\mathbf{A}[\mathbf{X}]}(\mathbf{M}) = n > 0\) . Considering as above an exact sequence \(0 \to \mathbf{N} \to \mathbf{L} \to \mathbf{M} \to 0\) , where \(\mathbf{L}\) is a free \(\mathbf{A}[\mathbf{X}]\) -module, we obtain an exact sequence of \(\mathbf{A}\) -modules
\end{enumerate}

\[
0 \rightarrow \mathrm{M} \rightarrow \mathrm{N} / \mathrm{XN} \rightarrow \mathrm{L} / \mathrm{XL} \rightarrow \mathrm{M} \rightarrow 0.
\]

By b), we have \(\mathrm{dp}_{\mathbf{A}}(\mathbf{N} / \mathbf{X}\mathbf{N}) \leqslant \mathrm{dp}_{\mathbf{A}[\mathbf{X}]}(\mathbf{N}) = \mathrm{dp}_{\mathbf{A}[\mathbf{X}]}(\mathbf{M}) - 1 = n - 1\) ; since \(\mathrm{dp}_{\mathbf{A}}(\mathbf{L} / \mathbf{XL}) = 0\) , we deduce from the preceding exact sequence, applying X, p.~135, cor. 2 twice, that \(\mathrm{dp}_{\mathbf{A}}(\mathbf{M}) \leqslant n - 1\) . But, by a), we have

\[
\mathrm{dp} _ {\mathbf {A}} (\mathbf {M}) \geqslant \mathrm{dp} _ {\mathbf {A} [ \mathbf {X} ]} (\mathbf {M}) - 1 = n - 1,
\]

whence c).

THEOREM 1. --- Suppose A is commutative. Then

\[
\mathrm{dh} (\mathbf {A} [ \mathbf {X} ]) = \mathrm{dh} (\mathbf {A}) + 1.
\]

For every A{[}X{]}-module M, we have (lemma 3)

\[
\mathrm{dp} _ {\mathbf {A} [ \mathbf {X} ]} (\mathbf {M}) \leqslant \mathrm{dp} _ {\mathbf {A}} (\mathbf {M}) + 1 \leqslant \mathrm{dh} (\mathbf {A}) + 1
\]

hence \(\mathrm{dh}(\mathbf{A}[\mathbf{X}]) \leqslant \mathrm{dh}(\mathbf{A}) + 1\) ; conversely, if \(\mathbf{M}\) is an A-module, let \(\overline{\mathbf{M}}\) be the \(\mathbf{A}[\mathbf{X}]\) -module obtained by endowing \(\mathbf{M}\) with the structure for which \(\mathbf{X}\mathbf{M} = 0\) , then (lemma 3)

\[
\mathrm{d} p _ {\mathrm{A}} (\mathrm{M}) = \mathrm{d} p _ {\mathrm{A} [ \mathrm{X} ]} (\overline {{\mathrm{M}}}) - 1 \leqslant \mathrm{d} h (\mathrm{A} [ \mathrm{X} ]) - 1,
\]

hence dh(A) ≤ dh(A{[}X{]}) - 1.

COROLLARY 1. --- Suppose A is commutative. We have

\[
\mathrm{dh} (\mathrm{A} [ \mathrm{X} _ {1}, \dots , \mathrm{X} _ {n} ]) = \mathrm{dh} (\mathrm{A}) + n.
\]

This follows from the theorem by induction on n.

COROLLARY 2. --- Let K be a commutative field (resp. a principal ideal ring * or a Dedekind ring * which is not a field). Then dh(K{[}X₁, \ldots, Xₙ{]}) is equal to n (resp. n + 1). This follows from the fact that dh(K) = 0 (resp. dh(K) = 1).

